\SourceLeaf{386}%
\noindent [so dass jetzt die zweiten Ableitungen der Coordinaten $x,y,z$ nach $t$ und $u$ allein durch $m,A,B,C$ und deren Ableitungen nach $t$ und $u$ ausgedrückt sind.]
\Needspace{16\baselineskip}\begin{center}[2.]\end{center}
\noindent\hspace*{1.2em} Die Formel
\[
\frac{\partial a'}{\partial u}=\frac{\partial a''}{\partial t}
\]
gibt uns folgende Gleichung:
\[
\begin{aligned}
&a\left(\frac{\partial\partial\log m}{\partial t^2}+\frac{\partial\partial\log m}{\partial u^2}\right)+\frac{m''}m\left(a'+\frac{\partial\partial x}{\partial u^2}\right)\\
&\quad+A''(aA'+bB'+cC')+A(a''A'+b''B'+c''C')+c'B'-b'C'\\
&\quad+A\left(a\frac{\partial\partial A}{\partial t\partial u}+b\frac{\partial\partial B}{\partial t\partial u}+c\frac{\partial\partial C}{\partial t\partial u}\right)\\
&\quad+c\frac{\partial\partial B}{\partial t^2}-b\frac{\partial\partial C}{\partial t^2}=0.
\end{aligned}
\]
Und eben so [liefern uns die Formeln
\[
\begin{aligned}
\frac{\partial b'}{\partial u}&=\frac{\partial b''}{\partial t}\\
\frac{\partial c'}{\partial u}&=\frac{\partial c''}{\partial t}]
\end{aligned}
\]
zwei andere ganz ähnliche [Gleichungen]; multiplicirt man sie respective mit $\alpha,\beta,\gamma$ und addirt, so erhält man
\[
\begin{aligned}
&mm\left(\frac{\partial\partial\log m}{\partial t^2}+\frac{\partial\partial\log m}{\partial u^2}\right)\\
&\quad+(aA'+bB'+cC')(\alpha A''+\beta B''+\gamma C'')+\Sigma\alpha(c'B'-b'C')\\
&\quad+\Sigma\alpha\left(c\frac{\partial\partial B}{\partial t^2}-b\frac{\partial\partial C}{\partial t^2}\right)=0.
\end{aligned}
\]
\par\indent Die beiden letzten Theile sind
\[
=\frac{\partial\Sigma\alpha(cB'-bC')}{\partial t}-\Sigma a''(cB'-bC').
\]
Nun aber ist [wegen 12]:
\begin{equation}\tag*{[23]}
\left\{\begin{aligned}
cB'-bC'&=A\theta\\
aC'-cA'&=B\theta\\
bA'-aB'&=C\theta,
\end{aligned}\right.
\end{equation}
[und hierin darf man, wiederum nach 12, setzen:]
